% Dyadic Hilbert-marked and pro-2 idempotent-power effectivity extension.
% Standalone section; integrated into fixed126.

\subsection{Hilbert-marked Fox coordinates and residual profinite-power effectivity}
\label{subsec:dyadic-hilbert-marked-power}

The dyadic Fox--Wu--Tate completion has two complementary native
interfaces.  The first identifies the mod-$2$ Fox quadratic form in
marked $G_{\mathbf Q_2}(2)$ coordinates with the classical explicit
quadratic Hilbert pairing.  The second evaluates the profinite
idempotent exponents appearing in the marked full-$G_{\mathbf Q_2}$
presentation by convergent, residue-sector-specific \emph{integral}
matrix power series.  Only finite jets are claimed to be polynomial;
a global polynomial formula for the whole profinite-power operation
is neither assumed nor needed.

\begin{theorem}[Hilbert-marked Fox comparison at $\mathbf Q_2$]
\label{thm:dyadic-fox-hilbert-marked}
Choose the Roe--Turturean marked Demu\v{s}kin normalization
\cite{RoeTurtureanGQ2}
\[
P=G_{\mathbf Q_2}(2)
 =\langle a,s,y\mid a^2s^4[s,y]\rangle_{\mathrm{pro}-2},
\]
in which the unramified character
$\nu_{\rm ur}:P\twoheadrightarrow\mathbb Z_2$
takes values
\begin{equation}\label{eq:q2-unramified-marking}
 \nu_{\rm ur}(a)=-2,\qquad
 \nu_{\rm ur}(s)=1,\qquad
 \nu_{\rm ur}(y)=0.
\end{equation}
Let $\chi_a,\chi_s,\chi_y$ be the mod-$2$ characters
dual to these generators, and let
$\kappa_u\in H^1(G_{\mathbf Q_2},\mathbf F_2)
=H^1(P,\mathbf F_2)$ denote the Kummer class
of $u\in\mathbf Q_2^\times$.
Then
\begin{equation}\label{eq:q2-kummer-fox-markings}
 \boxed{\quad
 \kappa_{-1}=\chi_a,\qquad
 \kappa_{5}=\chi_s,\qquad
 \kappa_2=\chi_y+\epsilon\chi_s
 \ \text{for a unique }\epsilon\in\mathbf F_2.
 \quad}
\end{equation}
Thus the marked Fox basis and the Kummer basis
$(-1,5,2)$ have the \emph{same} nonalternating cup matrix
\begin{equation}\label{eq:q2-hilbert-fox-matrix}
 B_{\mathbf Q_2}=
 \begin{pmatrix}
 1&0&0\\
 0&0&1\\
 0&1&0
 \end{pmatrix},
\end{equation}
where entry $1$ means Hilbert symbol $-1$.
In particular, every quadratic Hilbert-symbol pairing on
$\mathbf Q_2^\times/\mathbf Q_2^{\times2}$ is reconstructed
from the completed Fox quadratic symbol and the two canonical
markings $\theta_P\bmod4$ and $\nu_{\rm ur}\bmod2$.
No claim is made that the third generator-dual character
$\chi_y$ is \emph{literally} $\kappa_2$ without the
additional one-bit marking $\epsilon$.
\end{theorem}

\begin{proof}
The cyclotomic orientation values in
Proposition~\ref{prop:dyadic-q2-explicit-anchor}
are $(-1,1,-1/3)$ on $(a,s,y)$.
Since $-1/3\equiv1\bmod4$, the Fox--Wu identity
(Theorem~\ref{thm:dyadic-fox-wu-bockstein})
gives
\[
 \kappa_{-1}=w_P=(\theta_P-1)/2\bmod2=\chi_a.
\]
The extension $\mathbf Q_2(\sqrt5)/\mathbf Q_2$
is the unramified quadratic extension: for example,
$t^2-t-1$ is irreducible modulo $2$ and has
discriminant $5$.  Its Kummer character is
therefore the unique nontrivial unramified quadratic
character, namely $\nu_{\rm ur}\bmod2=\chi_s$
by \eqref{eq:q2-unramified-marking}.

Write $b=2^\beta v$ and $c=2^\gamma w$
with $v,w$ odd.  The explicit dyadic
Hilbert-symbol formula is
\begin{equation}\label{eq:q2-hilbert-explicit}
 (b,c)_2=(-1)^{
 \varepsilon(v)\varepsilon(w)
 +\beta\lambda(w)+\gamma\lambda(v)},
 \quad
 \varepsilon(v)=\frac{v-1}{2}\bmod2,\qquad
 \lambda(v)=\frac{v^2-1}{8}\bmod2,
\end{equation}
as in the standard local reciprocity computation
\cite{RoeTurtureanGQ2,NSW}.
For $v=-1,5,1$ respectively, the pairs
$(\varepsilon(v),\lambda(v))$ are
$(1,0),(0,1),(0,0)$.
Consequently the only nontrivial Hilbert
pairings among $(-1,5,2)$ are
\[
 (-1,-1)_2=-1,\qquad
 (5,2)_2=(2,5)_2=-1,
\]
giving \eqref{eq:q2-hilbert-fox-matrix}.
This equals the independent Fox cup matrix
of Proposition~\ref{prop:dyadic-q2-explicit-anchor}.

Write $\kappa_2=u\chi_a+v\chi_s+w\chi_y$.
Its pairing with $\kappa_{-1}=\chi_a$
vanishes, hence $u=0$ by the Fox cup matrix.
Its pairing with $\kappa_5=\chi_s$ is nonzero,
hence $w=1$.
Thus $\kappa_2=\chi_y+v\chi_s$,
which proves \eqref{eq:q2-kummer-fox-markings}
with $\epsilon=v$.
Adding $\epsilon\chi_s$ to $\chi_y$ leaves the
Fox cup matrix unchanged, as required.
\end{proof}

\begin{remark}[The remaining Hilbert-to-Fox bit is genuine marking data]
\label{rem:q2-hilbert-shear-boundary}
The mark $\epsilon$ in
\eqref{eq:q2-kummer-fox-markings}
depends on the chosen lift of the
unramified coordinate and the remaining
Fox generator, not on the abstract
bilinear form.  In the chosen generator
coordinates one can write
$\epsilon=\kappa_2(s)$.
Neither the Fox cup matrix nor the
mod-$4$ orientation determines that value.
This is a precisely isolated \emph{one-bit
coordinate issue}, not a failure of finite
Galois cohomology or local duality.
\end{remark}

\begin{theorem}[Residue-sector formula for the profinite $2$-power projection]
\label{thm:dyadic-omega2-power-formula}
Let $(A,\mathfrak m,k)$ be a complete Noetherian
local $\mathcal O_E$-algebra with finite residue
field $k$ of characteristic $2$, and let
$X\in\mathrm{GL}_r(A)$ reduce to
$\bar X\in\mathrm{GL}_r(k)$ of order
\[
 N=2^a m,\qquad m\text{ odd}.
\]
Let
$\omega_2\in\widehat{\mathbf Z}
=\prod_\ell\mathbf Z_\ell$ be the idempotent
whose $2$-adic component is $1$ and whose
other components are $0$.
Choose any nonnegative integer $e<N$ with
\begin{equation}\label{eq:dyadic-omega2-CRT}
 e\equiv1\pmod{2^a},\qquad
 e\equiv0\pmod m,
\end{equation}
where the first congruence is vacuous when $a=0$,
and set
\[
 c=\frac{1-e}{N}\in\mathbb Z_2,
 \qquad S=X^N-I_r\in M_r(\mathfrak m).
\]
Then the intrinsic profinite power is the
following explicit convergent integral series:
\begin{equation}\label{eq:dyadic-omega2-binomial-power}
 \boxed{\quad
 X^{\omega_2}
 =X^e\sum_{j\ge0}
   \binom cj (X^N-I_r)^j.
 \quad}
\end{equation}
For every $t\ge1$ it has the finite-stage formula
\begin{equation}\label{eq:dyadic-omega2-mod-t}
 X^{\omega_2}\equiv
 X^e\sum_{j=0}^{t-1}\binom cj S^j
 \pmod{\mathfrak m^t}.
\end{equation}
The binomial coefficients lie in $\mathbb Z_2$
and require no inversion of $2$.
The result is independent of the chosen
CRT representative $e$, commutes with
conjugation, and is natural under continuous
local coefficient base change preserving the
residual chart.

For a variation $H\in M_r(A)$, put
\begin{equation}\label{eq:dyadic-omega2-power-derivative}
 H_N=\sum_{\ell=0}^{N-1}
  X^\ell H X^{N-1-\ell},\qquad
 H_e=\sum_{\ell=0}^{e-1}
  X^\ell H X^{e-1-\ell}.
\end{equation}
Then the full matrix Jacobian of the profinite
power is represented by the convergent series
\begin{equation}\label{eq:dyadic-omega2-jacobian}
\begin{aligned}
 D(X^{\omega_2})[H]
 ={}&
 H_e\sum_{j\ge0}\binom cj S^j\\
 &+
 X^e\sum_{j\ge1}\binom cj
 \sum_{\ell=0}^{j-1}
  S^\ell H_N S^{j-1-\ell}.
\end{aligned}
\end{equation}
Modulo $\mathfrak m^t$, the first sum
stops at $j=t-1$ and the second at $j=t$.
Thus both the power operation and its
noncommutative first derivative have
terminating matrix formulas at every
specified finite residual precision.
\end{theorem}

\begin{proof}
The kernel of
$\mathrm{GL}_r(A)\to\mathrm{GL}_r(k)$
is pro-$2$ in the $\mathfrak m$-adic
topology.  Hence the closure of the
cyclic group generated by $X$ is a
procyclic profinite group, whose odd-order
factor has order $m$, while its remaining
factor is pro-$2$.  Write its canonical
commuting decomposition as
$X=X_{(2)}X_{\rm odd}$.
Then $X^N=X_{(2)}^N$ belongs to
$1+M_r(\mathfrak m)$.
The congruences
\eqref{eq:dyadic-omega2-CRT} give
$X_{\rm odd}^e=1$ and
$e+Nc=1$ in $\mathbf Z_2$.
Consequently
\[
 X^e(X^N)^c
 =X_{(2)}^{e+Nc}X_{\rm odd}^e
 =X_{(2)}=X^{\omega_2}.
\]
The usual binomial expression
\[
 (I_r+S)^c=\sum_{j\ge0}\binom cj S^j
\]
agrees with continuous $\mathbf Z_2$
exponentiation: it follows by
approximating $c$ with integers and
using $\mathfrak m$-adic continuity.
Since $S\in M_r(\mathfrak m)$,
the $j$th summand lies in
$M_r(\mathfrak m^j)$ and the
series converges, giving
\eqref{eq:dyadic-omega2-binomial-power}
and \eqref{eq:dyadic-omega2-mod-t}.
The equality with intrinsic
$X^{\omega_2}$ proves independence,
conjugation covariance, and naturality.

For the derivative, the standard
noncommutative power rule gives
$D(X^N)[H]=H_N$ and
$D(X^e)[H]=H_e$.
Differentiate the convergent binomial
series termwise, using
\[
 D(S^j)[H]=
 \sum_{\ell=0}^{j-1}
 S^\ell H_N S^{j-1-\ell},
\]
and apply the Leibniz rule for its
leading factor $X^e$.  This is
\eqref{eq:dyadic-omega2-jacobian}.
The derivative of the $j$th binomial
term has at least $j-1$ factors of
$S\in\mathfrak m$, which gives
the asserted exact truncation
order modulo $\mathfrak m^t$.
\end{proof}

\begin{example}[Removal of an odd residual eigenvalue]
\label{ex:dyadic-omega2-odd-teich}
Let $E/\mathbf Q_2$ be unramified quadratic,
so $\mathcal O_E$ contains a primitive
third root of unity $\zeta_3$.
For $A=\mathcal O_E[[z]]$ take the rank-one
transport $X=\zeta_3(1+z)$.
Its residual order is $N=3$.
Choose $e=0$ and $c=1/3\in\mathbb Z_2$.
The formula gives
\[
 X^{\omega_2}
 =\left(X^3\right)^{1/3}
 =\left((1+z)^3\right)^{1/3}
 =1+z.
\]
Thus the explicit integral binomial
calculus really removes the odd
Teichm\"uller component while retaining
the complete pro-$2$ deformation factor.
\end{example}

\begin{theorem}[Finite-jet full-$G_{\mathbf Q_2}$ marked word solver]
\label{thm:dyadic-q2-marked-word-jet-solver}
Assume the Roe--Turturean marked
profinite presentation of $G_{\mathbf Q_2}$
\cite{RoeTurtureanGQ2}, with generators
$\sigma,\tau,x_0,x_1$, tame relation
$\tau^\sigma=\tau^2$, one additional
marked relation $r_{\rm wild}=1$,
and the requirement that the normal
closure of $x_0,x_1$ be pro-$2$.
In the displayed marked word,
the idempotent powers are
\[
 \sigma_2=\sigma^{\omega_2},
 \qquad
 u_i=(x_i\tau)^{\omega_2}\quad(i=0,1),
\]
and all remaining operations are
finite products, inverses, conjugations,
and ordinary integral powers of
these marked words.

Fix a genuine residual representation
$\bar\rho:G_{\mathbf Q_2}\to\mathrm{GL}_r(k_E)$,
and choose integral lifts of its four
residual generator matrices.
Let
\[
 A_{\rm word}
 =\mathcal O_E[[z_{\sigma,ij},z_{\tau,ij},
                z_{0,ij},z_{1,ij}]]_{1\le i,j\le r}
\]
be the \emph{ambient} formally smooth
completed matrix-coordinate ring on
these four lifted generators.
Each occurrence of $\omega_2$ on
$A_{\rm word}$ has the
convergent expression
\eqref{eq:dyadic-omega2-binomial-power}.
For every requested integer $t\ge1$:
\begin{enumerate}[label=\textnormal{(\roman*)}]
\item the full tame and marked wild
matrix relation maps are evaluated
modulo $\mathfrak m^t$ by a finite
sequence of matrix products and
the truncated binomial formula
\eqref{eq:dyadic-omega2-mod-t};
\item the entire first matrix Jacobian
of both marked relations is computed
modulo $\mathfrak m^t$ by the
finite derivative formula
\eqref{eq:dyadic-omega2-jacobian},
product rules, and inverse-matrix
differentiation;
\item the pro-$2$ marked-normal-closure
condition imposes no additional
equation on any Artin lift with
this genuine residual image;
\item the quotient of the ambient
$A_{\rm word}$ by the scalar entries
of the two full convergent relation
matrices is the completed framed
deformation ring $R_{\bar\rho}^{\square}$;
\item at the first-order residual
thickening the kernel of the
\emph{ambient} marked matrix
Jacobian is canonically the framed
cocycle space $Z^1(G_{\mathbf Q_2},
\operatorname{ad}\bar\rho)$.
Quotienting by infinitesimal
conjugation gives
$H^1(G_{\mathbf Q_2},
\operatorname{ad}\bar\rho)$.
\end{enumerate}
This is a \emph{terminating
formal-native algorithm} at every
prescribed finite order.  It does
not identify the naive two-relation
Jacobian cokernel with $H^2$ or
the naive two-relator Koszul complex
with an integral group-resolution.
The independent lci minimal-sequence
and intrinsic derived-cochain
constructions of
Theorem~\ref{thm:word-lci-reduction}
and Theorem~\ref{thm:family-sylow-fox}
remain necessary for those statements.
\end{theorem}

\begin{proof}
For each of the residual matrices of
$\sigma,x_0\tau,x_1\tau$ its order in
$\mathrm{GL}_r(k_E)$ is finite.
Apply Theorem~\ref{thm:dyadic-omega2-power-formula}
to these three matrices with their
own CRT exponents.  Every occurrence
of a profinite idempotent power in
the Roe--Turturean relation becomes
an integral convergent analytic
matrix expression on the completed
residual chart.
Since the remaining word operations
are finite, truncating each
constituent to the requested
$\mathfrak m$-adic order evaluates
the relation matrices by finitely
many operations.  The explicit
Jacobian in
\eqref{eq:dyadic-omega2-jacobian}
and the ordinary matrix product,
inverse, and conjugation rules
prove (i)--(ii).

The marked wild normal subgroup of
$G_{\mathbf Q_2}$ is pro-$2$,
so its image under the genuine
residual representation is a
finite $2$-group.  On every Artin
local lift, the reduction kernel
is a successive extension of
additive groups in characteristic
$2$, hence a $2$-group.
The normal closure of the lifted
wild image is therefore an
extension of finite $2$-groups,
so is again a $2$-group.
This is precisely the residual
automaticity theorem
(Theorem~\ref{thm:marked-residual-strictification}),
giving (iii).  Because the Roe--Turturean
presentation is universal among the
marked finite quotients, the same
residual automaticity identifies the
functor of solutions to the completed
two-relation equations in $A_{\rm word}$
with the framed deformation functor.
This proves (iv).

On the dual-number thickening
$k_E[\varepsilon]/(\varepsilon^2)$,
every lift is
$\rho_\varepsilon(g)
=(1+\varepsilon c(g))\bar\rho(g)$.
The defining representation
equation is the continuous
crossed-cocycle equation
$c(gh)=c(g)+g\cdot c(h)$.
By the marked presentation and
(iii), it is equivalent to the
linearization of the two
displayed marked relations,
computed by (ii).  This identifies
the tangent kernel with $Z^1$.
Infinitesimal conjugation is
exactly the coboundary action,
giving (v).
The last assertion follows from
the distinction between formal
relation equations and a complete
projective group resolution
already made in
Warning~\ref{warn:q2-marked-not-resolution}.
\end{proof}

\begin{warning}[Finite residual jets are not a globally polynomial presentation]
\label{warn:dyadic-marked-word-not-algebraic-polynomial}
The matrices in
\eqref{eq:dyadic-omega2-binomial-power}
are generally infinite convergent
power series on one chosen completed
residual chart.  Their reductions
at each finite order are literal
finite matrix expressions, but
a single finite polynomial in
unrestricted generic matrices
does not follow.
The exact native output of
Theorem~\ref{thm:dyadic-q2-marked-word-jet-solver}
is a recursively effective
formal marked-word atlas, with
strict integral tangent readout.
Global finite-type algebraization,
a uniform marked presentation for
all dyadic $K$, and a preferred
raw Herr-to-marked-word chain
quasi-isomorphism remain stronger
separate questions.  None is
needed by the established
all-Sylow--Cayley--Hamilton
derived faithful atlas.
\end{warning}
